\section{Implementação, Profiling e Auditoria}
\label{sec:implementacao}

O sistema experimental foi implementado em Python~3 sob uma arquitetura modular que separa o
carregamento de instâncias, os núcleos algorítmicos, o orquestrador de experimentos e os
módulos de análise estatística. As estruturas numéricas densas apoiam-se em arranjos contíguos
em memória de precisão dupla, reservando o uso de bibliotecas de grafos estritamente à etapa de
emparelhamento perfeito de peso mínimo (algoritmo Blossom) e oferecendo suporte opcional a
processamento paralelo em GPU via CUDA para as versões vetorizadas da busca exata. O
desenvolvimento contou com apoio controlado de ferramentas de inteligência artificial generativa
para prototipação, auditoria e revisão; o registro consolidado de prompts do grupo encontra-se
documentado separadamente no repositório do projeto.

\subsection{Construção e Auditoria da Base de Instâncias (\texttt{EUC\_2D})}
\label{sec:base_dados}

Para que a comparação entre heurísticas geométricas (como a inserção a partir do fecho convexo),
algoritmos métricos e busca exata fosse matematicamente consistente, o conjunto de instâncias
extraído da biblioteca clássica TSPLIB95 \cite{reinelt1991tsplib} passou por quatro decisões
metodológicas e verificações de integridade:

\begin{enumerate}
  \item \textbf{Uso estrito de coordenadas bidimensionais nativas}: Uma inspeção direta dos
  arquivos oficiais da TSPLIB95 constatou que diversas instâncias pequenas sugeridas
  no enunciado do projeto (como \texttt{gr24}, \texttt{fri26} e \texttt{swiss42}) fornecem apenas matrizes
  explícitas de distâncias entre cidades, sem coordenadas cartesianas $(x, y)$ associadas.
  Projetar essas matrizes no plano por técnicas de imersão multidimensional introduziria
  distorções métricas e posições artificiais, comprometendo tanto o cálculo da envoltória
  convexa quanto as métricas de distribuição espacial. Por essa razão, descartamos qualquer
  instância sem coordenadas bidimensionais nativas. Também removemos \texttt{dantzig42}
  ($N = 42$) por situar-se fora dos limites adotados para a divisão de escalas entre a faixa
  principal ($N \le 40$) e a faixa de estresse ($N \in [45, 75]$) definidas no escopo do trabalho.

  \item \textbf{Balanceamento em 15 instâncias (80\% faixa principal e 20\% faixa de estresse)}:
  Incluímos integralmente todas as instâncias oficiais da TSPLIB95 com coordenadas 2D nativas
  na faixa $N \le 40$ (\texttt{burma14}, \texttt{ulysses16}, \texttt{ulysses22} e
  \texttt{bays29}, esta última a partir de sua seção oficial de coordenadas de exibição), com
  verificação byte a byte em relação aos arquivos originais. Para preencher de maneira regular
  o intervalo $N \in [10, 40]$ necessário às regressões de complexidade, construímos recortes
  prefixados determinísticos (os $k$ primeiros pontos nativos) das instâncias euclidianas
  clássicas \texttt{eil51}, \texttt{berlin52} e \texttt{st70}. Essa composição resulta em
  exatamente 12 instâncias ($80\%$) na faixa principal (\texttt{eil51\_n10},
  \texttt{berlin52\_n12}, \texttt{burma14}, \texttt{ulysses16}, \texttt{eil51\_n18},
  \texttt{berlin52\_n20}, \texttt{ulysses22}, \texttt{eil51\_n25}, \texttt{bays29},
  \texttt{berlin52\_n33}, \texttt{eil51\_n38} e \texttt{st70\_n40}) e 3 instâncias completas
  ($20\%$) na faixa de estresse (\texttt{att48}, \texttt{berlin52} e \texttt{st70}), cobrindo o
  intervalo $N \in [10, 70]$.

  \item \textbf{Distância euclidiana contínua versus arredondamento inteiro}: Na especificação
  histórica da TSPLIB, pesos \texttt{EUC\_2D} costumam ser arredondados para o inteiro mais
  próximo. Todavia, instâncias como \texttt{burma14}, \texttt{ulysses16} e \texttt{ulysses22}
  expressam suas coordenadas nativas em graus decimais, produzindo distâncias euclidianas muito
  pequenas entre pontos próximos (frequentemente inferiores a uma unidade). Aplicar
  arredondamento inteiro sobre essas coordenadas colapsaria arestas curtas para custo zero e
  criaria dezenas de empates artificiais, descaracterizando a geometria do grafo. Optou-se,
  portanto, por avaliar $c(u, v) = \sqrt{(x_u - x_v)^2 + (y_u - y_v)^2}$ de forma contínua em
  precisão dupla (\texttt{float64}), mantendo o modo inteiro apenas como opção parametrizável
  para testes controlados de grafos com pesos discretos.

  \item \textbf{Certificação de um gabarito ótimo estritamente \texttt{EUC\_2D}}: Ao auditar
  os arquivos de rotas ótimas distribuídos com a TSPLIB, detectou-se uma incompatibilidade
  métrica relevante. Como \texttt{bays29} teve sua rota oficial calculada originalmente sobre
  uma matriz rodoviária explícita, essa rota não é necessariamente ótima sob a distância
  euclidiana reta entre suas coordenadas; situação análoga ocorre em \texttt{ulysses16} e
  \texttt{ulysses22} (originalmente definidos sob métrica geográfica esférica) e em recortes
  prefixados que não possuem arquivo oficial. Caso tais rotas fossem usadas como referência de
  ótimo, heurísticas euclidianas poderiam obter custos inferiores ao suposto ótimo, gerando
  erros relativos negativos sem significado matemático. Para garantir rigor na medição da
  qualidade das aproximações e na análise de sensibilidade (inclusive para instâncias em que o
  \textit{Branch-and-Bound} atinge o limite de tempo), construímos um oráculo estático 
  baseado na formulação linear inteira de Dantzig--Fulkerson--Johnson \cite{dantzig1954solution}
  com inserção iterativa de desigualdades de eliminação de subciclos
  ($\sum_{i, j \in S} x_{ij} \le |S| - 1$ para todo subconjunto desconexo $S \subsetneq V$).
  Esse procedimento certifica o valor ótimo exato contínuo $OPT$ de todas as 15 instâncias e
  atua exclusivamente como gabarito estático para o cálculo do erro percentual, nunca
  substituindo o \textit{Branch-and-Bound} na avaliação do algoritmo exato.
\end{enumerate}
